\documentclass[10pt,a4paper]{amsart}
\usepackage{amsmath,amssymb,amsthm,graphicx}
\usepackage{verbatim}
\newenvironment{reference}{\comment}{\endcomment}
%\newenvironment{reference}{}{}
\newenvironment{slogan}{\comment}{\endcomment}
\newenvironment{history}{\comment}{\endcomment}

% For commutative diagrams we use Xy-pic
\usepackage[all]{xy}

% We use 2cell for 2-commutative diagrams.
\xyoption{2cell}
% Theorem environments.
%
\theoremstyle{plain}
\newtheorem{theorem}[subsection]{Theorem}
\newtheorem{proposition}[subsection]{Proposition}
\newtheorem{lemma}[subsection]{Lemma}

\theoremstyle{definition}
\newtheorem{definition}[subsection]{Definition}
\newtheorem{example}[subsection]{Example}
\newtheorem{exercise}[subsection]{Exercise}
\newtheorem{situation}[subsection]{Situation}

\theoremstyle{remark}
\newtheorem{remark}[subsection]{Remark}
\newtheorem{remarks}[subsection]{Remarks}

\numberwithin{equation}{subsection}

% Macros
%
\def\lim{\mathop{\mathrm{lim}}\nolimits}
\def\colim{\mathop{\mathrm{colim}}\nolimits}
\def\Spec{\mathop{\mathrm{Spec}}}
\def\Hom{\mathop{\mathrm{Hom}}\nolimits}
\def\Ext{\mathop{\mathrm{Ext}}\nolimits}
\def\SheafHom{\mathop{\mathcal{H}\!\mathit{om}}\nolimits}
\def\SheafExt{\mathop{\mathcal{E}\!\mathit{xt}}\nolimits}
\def\Sch{\mathit{Sch}}
\def\Mor{\mathop{\mathrm{Mor}}\nolimits}
\def\Ob{\mathop{\mathrm{Ob}}\nolimits}
\def\Sh{\mathop{\mathit{Sh}}\nolimits}
\def\NL{\mathop{N\!L}\nolimits}
\def\CH{\mathop{\mathrm{CH}}\nolimits}
\def\proetale{{pro\text{-}\acute{e}tale}}
\def\etale{{\acute{e}tale}}
\def\QCoh{\mathit{QCoh}}
\def\Ker{\mathop{\mathrm{Ker}}}
\def\Im{\mathop{\mathrm{Im}}}
\def\Coker{\mathop{\mathrm{Coker}}}
\def\Coim{\mathop{\mathrm{Coim}}}

% Boxtimes
%
\DeclareMathSymbol{\boxtimes}{\mathbin}{AMSa}{"02}

%
% Macros for moduli stacks/spaces
%
\def\QCohstack{\mathcal{QC}\!\mathit{oh}}
\def\Cohstack{\mathcal{C}\!\mathit{oh}}
\def\Spacesstack{\mathcal{S}\!\mathit{paces}}
\def\Quotfunctor{\mathrm{Quot}}
\def\Hilbfunctor{\mathrm{Hilb}}
\def\Curvesstack{\mathcal{C}\!\mathit{urves}}
\def\Polarizedstack{\mathcal{P}\!\mathit{olarized}}
\def\Complexesstack{\mathcal{C}\!\mathit{omplexes}}
% \Pic is the operator that assigns to X its picard group, usage \Pic(X)
% \Picardstack_{X/B} denotes the Picard stack of X over B
% \Picardfunctor_{X/B} denotes the Picard functor of X over B
\def\Pic{\mathop{\mathrm{Pic}}\nolimits}
\def\Picardstack{\mathcal{P}\!\mathit{ic}}
\def\Picardfunctor{\mathrm{Pic}}
\def\Deformationcategory{\mathcal{D}\!\mathit{ef}}
\usepackage{cleveref}
\crefname{originallemma}{original lemma}{original lemmas}
\Crefname{originallemma}{Original lemma}{Original lemmas}
\makeatletter
\newlabel{lemma-map-between-finite}{{05GE}{source}{}{sourcepointer-1}{}}
\newlabel{lemma-map-between-finite@cref}{{[originallemma][0][]05GE}{[1][0][]source}{}{}{}}
\makeatother
\title{The locus where a map from a finite module to a finitely presented module is an isomorphism is open}
\author{The Stacks Project Authors}
\date{stacks-project:algebra.tex:local-snapshot-20261012; GNU FDL 1.2 or later}
\begin{document}
\maketitle
\noindent\textit{Editorial note.} All delivered mathematical text is the author's own, copied from the pinned Stacks Project chapter file algebra.tex; no statement or proof step was rewritten, reordered or removed.\par\medskip
\noindent\textit{Editorial note.} Preamble provenance, as required for this source: the frozen snapshot stages the chapter's own preamble file beside it (preamble.tex, 217 lines, 8116 bytes, the same shared Stacks-wide preamble the published chow and etale-cohomology cards use); the first line of the pinned chapter file is \textbackslash{}input\{preamble\}, so no preamble text lives inside the chapter itself. The card takes that file's lines 8-18 (the verbatim package and the three \textbackslash{}comment-based environments reference/slogan/history at lines 9-12, section 13.2, together with the Xy-pic setup) and its lines 154 to the runtime end of the file (the native theorem declarations and mathematical macros). Not one character was changed, no package was added, no macro was defined, and the excerpt serves only as preamble, never as a span coordinate.\par\medskip
\noindent\textit{Editorial note.} The same-chapter pointers are located at their \textbackslash{}label lines inside the pinned chapter file; the cross-chapter pointers are located in the project's own tag list staged in the snapshot (sources/stacks-chow/tags.txt), at the verbatim lines carrying their canonical tags.\par\medskip
\noindent\textit{Editorial note.} The counted proof refers to same-chapter statements by label; they are recorded as navigation stubs pointed at their \textbackslash{}label lines, and the arguments consuming them stay inside the counted span.\par\medskip
\noindent\textit{Editorial note.} The counted spans contain no \textbackslash{}cite, so the card carries no bibliography.\par\medskip
\noindent\textit{Allowed standard prerequisites (editorial).}\begin{itemize}
\item A finite presentation of the target module. Used at the start of the counted proof (lines 19638-19640). Lets preimages of the generators and of the relations be written down explicitly.
\item Finitely many conditions spread out to a basic open. Used at the end of the counted proof (lines 19650-19660). Produces elements whose simultaneous nonvanishing defines an open subset of the isomorphism locus.
\end{itemize}\medskip
\ifdefined\hypertarget\hypertarget{sourcepointer-1}{}\fi
\noindent\textit{Source reference.} Original lemma 05GE: sources/stacks-algebra/pinned-original-source.tex, source lines 19601--19601. Navigation pointer for \textbackslash{}ref\{lemma-map-between-finite\} used inside the counted span (line 19656); the argument that consumes it is retained in the span.\par\medskip
\begin{lemma}
\label{lemma-map-between-finitely-presented}
Let $R$ be a ring. Let $\varphi : M \to N$ be a map of $R$-modules
with $M$ finite and $N$ finitely presented. Then
$$
U = \{\mathfrak p \subset R \mid
\varphi_{\mathfrak p} : M_{\mathfrak p} \to N_{\mathfrak p}
\text{ is an isomorphism}\}
$$
is an open subset of $\Spec(R)$.
\end{lemma}

\begin{proof}
Let $\mathfrak p \in U$. Pick a presentation
$N = R^{\oplus n}/\sum_{j = 1, \ldots, m} R k_j$.
Denote $e_i$ the image in $N$ of the $i$th basis vector of $R^{\oplus n}$.
For each $i \in \{1, \ldots, n\}$ choose an element
$m_i \in M_{\mathfrak p}$ such that $\varphi(m_i) = f_i e_i$
for some $f_i \in R$, $f_i \not \in \mathfrak p$. This is possible
as $\varphi_{\mathfrak p}$ is an isomorphism. Set $f = f_1 \ldots f_n$
and let $\psi : R_f^{\oplus n} \to M_f$ be the map which maps
the $i$th basis vector to $m_i/f_i$. Note that
$\varphi_f \circ \psi$ is the localization
at $f$ of the given map $R^{\oplus n} \to N$.
As $\varphi_{\mathfrak p}$ is an isomorphism we
see that $\psi(k_j)$ is an element of $M$ which maps to zero
in $M_{\mathfrak p}$. Hence we see that there exist
$g_j \in R$, $g_j \not \in \mathfrak p$ such that $g_j \psi(k_j) = 0$.
Setting $g = g_1 \ldots g_m$, we see that $\psi_g$ factors through
$N_{fg}$ to give a map $\chi : N_{fg} \to M_{fg}$. By construction
$\chi$ is a right inverse to $\varphi_{fg}$. It follows that
$\chi_\mathfrak p$ is an isomorphism. By Lemma \ref{lemma-map-between-finite}
there is an $h \in R$, $h \not \in \mathfrak p$ such that
$\chi_h : N_{fgh} \to M_{fgh}$ is surjective.
Hence $\varphi_{fgh}$ and $\chi_h$ are mutually inverse maps,
which implies that $D(fgh) \subset U$ as desired.
\end{proof}

\end{document}
